\section{Probability - The Science of Uncertainty and Data (6.431x)}
\subsection{Probability models and axioms}

A probabilistic model describes the behavior of a random experiment, such as flipping a coin, rolling dice, or throwing a dart. Building such a model has two steps: first, we list all possible outcomes; second, we assign probabilities to them. In this section, we focus on the first step — describing the possible outcomes clearly and systematically.

\subsubsection{Sample space}

All possible outcomes of an experiment form a set, denoted by \(\Omega\), called the sample space. Each element of \(\Omega\) represents one possible result. The outcomes must be mutually exclusive (only one can happen at a time) and collectively exhaustive (some outcome from \(\Omega\) must always happen).

How detailed our sample space is depends on what we care about in the experiment. It should include all relevant distinctions but ignore unimportant ones. Choosing a good sample space is part of the modeling process.

\bigskip
If we flip a coin, a simple sample space is
\[
\Omega = \{\text{heads}, \text{tails}\}.
\]
This covers all outcomes and they don’t overlap.

If we also record whether it’s raining, we could define a larger sample space,
\[
\Omega' = \{\text{heads and rain}, \text{heads and no rain}, \text{tails and rain}, \text{tails and no rain}\}.
\]
This is valid too, but since weather usually doesn’t affect a coin toss, it’s an unnecessary detail. We include such extra factors only when they matter for the problem.

\subsubsection{Discrete and continuous sample spaces}

Sample spaces can be discrete or continuous.

For example, rolling a four-sided die twice and recording the two results gives
\[
\Omega = \{(i,j): i,j \in \{1,2,3,4\}\},
\]
with 16 possible ordered pairs like \((1,1)\) or \((3,4)\). When experiments have several stages, tree diagrams can help visualize the structure of outcomes.

Other sample spaces are continuous. For instance, when throwing a dart at a unit square, each landing point is a pair of coordinates \((x,y)\):
\[
\Omega = \{(x,y): 0 \le x \le 1,\ 0 \le y \le 1\}.
\]
Here, every point inside the square is a possible outcome.

\subsubsection{Events and probability laws}

Once we have a sample space, we can talk about events. An event \(A\) is any subset of \(\Omega\). After an experiment, if the observed outcome \(\omega\) belongs to \(A\), we say the event occurred.

In discrete cases, we can assign probabilities to single outcomes. In continuous cases, each exact point has probability 0, so we assign probabilities to regions instead (for example, “the dart lands in the upper half of the square”). The probability law gives each event \(A\) a value \(P(A)\), showing how likely that event is.

\subsubsection{Probability axioms and basic set notation}

Probabilities follow three basic rules, called axioms:

\begin{itemize}
    \item \textbf{Nonnegativity}: For any event \(A\), \(P(A) \ge 0\).
    \item \textbf{Normalization}: \(P(\Omega) = 1\), since something in the sample space must occur.
    \item \textbf{Additivity}: If two events \(A\) and \(B\) are disjoint (that is, \(A \cap B = \emptyset\)), then
    \[
    P(A \cup B) = P(A) + P(B).
    \]
\end{itemize}

You can think of probability as a kind of “mass” spread over \(\Omega\). If two events don’t overlap, the total mass on their union is just the sum of their individual masses.

\newpage
For two events \(A\) and \(B\):
\begin{itemize}
    \item The intersection \(A \cap B\) is the set of outcomes that belong to both \(A\) and \(B\).
    \item The union \(A \cup B\) is the set of outcomes that belong to either \(A\) or \(B\) (or both).
    \item Two events are disjoint if \(A \cap B = \emptyset\), meaning they have no outcomes in common.
\end{itemize}


\subsubsection{Simple properties derived from the axioms}

The three axioms look simple, but they already imply many useful and intuitive properties. For example, we can prove that probabilities are always at most 1, and that the probability of the empty set is 0, instead of assuming these as extra axioms.

Consider an event \(A\) and its complement \(A^c\), which contains all outcomes in \(\Omega\) that are not in \(A\). The set \(A\) and its complement \(A^c\) are disjoint, and their union is the entire sample space:
\[
A \cup A^c = \Omega, \qquad A \cap A^c = \emptyset.
\]
Using normalization and additivity, we get
\[
P(\Omega) = P(A \cup A^c) = P(A) + P(A^c) = 1.
\]
From this relation we also obtain
\[
P(A) = 1 - P(A^c).
\]
Because probabilities are nonnegative, \(P(A^c) \ge 0\), so \(P(A) \le 1\). This shows that probabilities are always bounded above by 1, without adding a separate axiom.

If we apply the same relation to the special case where \(A = \Omega\), then the complement \(A^c\) is the empty set \(\emptyset\). We already know that \(P(\Omega) = 1\), so
\[
1 = P(\Omega) = P(\Omega) + P(\emptyset),
\]
which forces \(P(\emptyset) = 0\). This matches the idea that the empty set is “impossible,” so its probability should be zero.

Finally, the additivity axiom for two disjoint events extends to any finite collection of disjoint events. If \(A_1, A_2, \dots, A_k\) are pairwise disjoint, then
\[
P\Big(\bigcup_{i=1}^k A_i\Big) = \sum_{i=1}^k P(A_i).
\]
A common special case is when we have a finite subset of the sample space consisting of points \(s_1, \dots, s_k\). We can view this set as the union of the singletons \(\{s_1\}, \dots, \{s_k\}\), which are disjoint. Then the probability of the finite set is just the sum of the probabilities of its individual elements. In practice we often write \(P(s_i)\) instead of \(P(\{s_i\})\) and talk informally about “the probability of an outcome \(s_i\)” even though probabilities are formally defined on sets.

\bigskip

The axioms also give us some useful monotonicity and union properties. Suppose we have two events \(A\) and \(B\) with \(A \subseteq B\). Intuitively, \(B\) is “larger,” so its probability should be at least as big as that of \(A\). Formally, we can write
\[
B = A \cup \big(B \cap A^c\big),
\]
where \(A\) and \(B \cap A^c\) are disjoint. By additivity,
\[
P(B) = P(A) + P\big(B \cap A^c\big) \ge P(A),
\]
since probabilities are nonnegative.

We can also express the probability of a union when the sets are not disjoint. If \(A\) and \(B\) overlap, then the union \(A \cup B\) can be decomposed into three disjoint parts: the part of \(A\) outside \(B\), the intersection \(A \cap B\), and the part of \(B\) outside \(A\). If we call their probabilities \(a\), \(b\), and \(c\), then
\[
P(A \cup B) = a + b + c.
\]
On the other hand,
\[
P(A) = a + b, \qquad P(B) = b + c, \qquad P(A \cap B) = b.
\]
Putting these together gives the standard identity
\[
P(A \cup B) = P(A) + P(B) - P(A \cap B).
\]

A direct corollary is the union bound. Since \(P(A \cap B) \ge 0\), we get
\[
P(A \cup B) \le P(A) + P(B),
\]
which is useful when we only need an upper bound on a probability.

The same idea extends to three sets. For events \(A\), \(B\), and \(C\), we can rewrite their union as a union of three disjoint pieces: \(A\) itself, the part of \(B\) outside \(A\), and the part of \(C\) outside both \(A\) and \(B\). Using additivity on these disjoint parts gives a formula for \(P(A \cup B \cup C)\) in terms of simpler probabilities. More generally, we can keep breaking unions into disjoint pieces to get expressions for probabilities of unions of several events.

\subsubsection{A discrete example and the discrete uniform law}

Consider again two rolls of a tetrahedral die. The sample space has 16 outcomes, each an ordered pair \((X,Y)\), where \(X\) is the first roll and \(Y\) is the second. Assume all 16 outcomes are equally likely, so each has probability \(1/16\).

As a first example, look at the event \(\{X = 1\}\): the first roll is 1. In the grid of all \((X,Y)\) pairs, this event contains the four outcomes where the first coordinate is 1 and the second is 1, 2, 3, or 4. Since each outcome has probability \(1/16\), we get
\[
P(X = 1) = 4 \cdot \frac{1}{16} = \frac{1}{4}.
\]

Now define \(Z\) to be the smaller of the two numbers in the two rolls. For example, if \(X = 2\) and \(Y = 3\), then \(Z = 2\). For \(P(Z = 4)\), this happens only if both \(X\) and \(Y\) are 4, so there is exactly one outcome in the event and
\[
P(Z = 4) = \frac{1}{16}.
\]
For \(\{Z = 2\}\), we need one roll equal to 2 and the other at least 2. This event contains \((2,2)\), \((2,3)\), \((2,4)\), \((3,2)\), and \((4,2)\), for a total of five outcomes. Thus
\[
P(Z = 2) = 5 \cdot \frac{1}{16}.
\]

This is an instance of the discrete uniform law. If the sample space \(\Omega\) is finite with \(n\) elements and all are equally likely, then, since \(P(\Omega) = 1\), we must have
\[
P(\{\omega\}) = \frac{1}{n}
\]
for each outcome \(\omega\). If an event \(A\) has exactly \(k\) elements, then
\[
P(A) = k \cdot \frac{1}{n} = \frac{k}{n}.
\]
So in a discrete uniform model, computing probabilities amounts to counting: we count all possible outcomes \(n\), and the favorable ones \(k\), then take the ratio \(k/n\).

\subsubsection{A continuous example}

Now consider a probability calculation in a continuous sample space. The experiment is again throwing a dart at the unit square, and the outcome is the point \((x,y)\) where the dart lands. The sample space is the unit square, and we assume a uniform law: the probability of any region in the square is equal to its area. This is natural if we think the dart is equally likely to land anywhere.

Take the event that the sum of the coordinates is at most \(1/2\):
\[
A = \{(x,y) : x + y \le 1/2\}.
\]
In the unit square this is a right triangle under the line \(x + y = 1/2\). Since probability equals area,
\[
P(A) = \frac{1}{2} \cdot \frac{1}{2} \cdot \frac{1}{2} = \frac{1}{8}.
\]
As another example, consider the event that the dart hits exactly the point \((0.5, 0.3)\). This event is a single point with area 0, so its probability is 0. The same holds for any single point in the unit square: in a continuous model, individual points have probability 0, while regions with positive area can have positive probability.

This example fits a general four-step recipe for probability problems. Starting from a word description of an experiment, we:
\begin{enumerate}
    \item Specify the sample space.
    \item Choose a probability law (a model for how likely different outcomes are).
    \item Describe the event of interest precisely, ideally also with a picture.
    \item Use the probability law and the event description to compute the desired probability.
\end{enumerate}
Sometimes the law is given explicitly for all events; in other cases we only know it for some basic events and need extra work to obtain the probability we want.

\subsubsection{Countable additivity and infinite sample spaces}

We now look at a discrete but infinite sample space. Imagine an experiment whose outcome is a positive integer, for example the number of coin tosses until the first heads appears. The sample space is \(\{1,2,3,\dots\}\). Suppose we are told that
\[
P(\{n\}) = \frac{1}{2^n}, \quad n = 1,2,3,\dots
\]
First we check that this is a valid probability law by summing over all \(n\):
\[
\sum_{n=1}^\infty \frac{1}{2^n}
= \frac{1}{2} \sum_{n=0}^\infty \Big(\frac{1}{2}\Big)^n
= \frac{1}{2} \cdot \frac{1}{1 - 1/2} = 1.
\]

Now consider the event that the outcome is even, the infinite set \(\{2,4,6,\dots\}\). We can write it as a union of singletons:
\[
\{2,4,6,\dots\} = \{2\} \cup \{4\} \cup \{6\} \cup \dots
\]
These sets are disjoint, so we would like to write
\[
P(\text{even}) = P(\{2\}) + P(\{4\}) + P(\{6\}) + \dots
\]
and evaluate the resulting infinite series. To make this step valid in general, we extend the additivity axiom from finite unions to countably infinite unions.

The extra axiom, called countable additivity, says: if \(A_1,A_2,A_3,\dots\) is a sequence of pairwise disjoint events, then
\[
P\Big(\bigcup_{n=1}^\infty A_n\Big) = \sum_{n=1}^\infty P(A_n).
\]
The word “sequence” matters: the events can be listed in order \(A_1,A_2,\dots\), so we can form a countable sum. This axiom lets us compute probabilities of many events in infinite sample spaces by summing over simpler disjoint pieces.

It is also important to see what countable additivity does not say. In the continuous unit square model, the square is the union of singletons (one per point), and each singleton has probability 0. If we tried to add their probabilities, we would get a sum of zeros, but the probability of the whole square should be 1. The catch is that the unit square is uncountable: its points cannot be arranged as a sequence \(A_1,A_2,\dots\), so the axiom does not apply. In short, some infinite sets are countable (like the integers) and others are uncountable (like the real line or the unit square), and countable additivity only covers the countable case.

There are deeper measure-theoretic issues in the background, especially when we try to define probability on all subsets of a continuous space. For this course, the key point is that standard constructions, like defining probability on the unit square via area, can be made rigorous and do satisfy countable additivity for all “nice” sets we will use. The strange sets where problems arise are mostly of theoretical interest and do not appear in typical applications.

\subsubsection{Interpretations and uses of probability}

To close, it is useful to think about what probability really means. In the narrow sense, probability theory is a branch of mathematics: we start from axioms, study models that satisfy them, and prove theorems. This view does not fix any particular real-world interpretation.

One common interpretation is in terms of long-run frequencies. For example, if we have a fair coin and toss it many times, the fraction of heads is expected to approach \(1/2\). In this view, the probability of an event is the limiting frequency with which it occurs in a long sequence of repeated trials. This fits well for repeatable experiments like coin tosses.

But not all probability statements are like that. A sentence such as “the current president of my country will be reelected with probability \(0.7\)” does not describe an experiment we can repeat infinitely many times. In such cases it is more natural to see probabilities as measuring degrees of belief: they summarize how strongly we believe an event will happen. Thinking in terms of bets, a probability expresses how much we are willing to stake on an outcome.

This raises the worry that probabilities are just subjective. The answer is that probability theory at least gives consistent rules for reasoning under uncertainty. If our probability model is a reasonable description of the situation, those rules help us make predictions and decisions in a systematic way.

Whether our decisions are actually good depends on whether we chose a good model. Building such models from data is the task of statistics and statistical inference. The real world generates data; statistics uses these data to construct probabilistic models; and probability theory then analyzes these models to produce predictions or decisions that we can apply back to the real world.


\subsubsection{Sets, complement, empty set, and subsets}

A set is a collection of distinct elements, usually denoted by a capital letter like \(S\). We can specify a set by listing its elements inside braces, for example \(\{a,b,c,d\}\), which is a finite set. Other sets, such as the set of all real numbers, are infinite. We write \(x \in S\) if \(x\) is an element of \(S\), and \(x \notin S\) if \(x\) does not belong to \(S\).

Sometimes it is convenient to define a set by a property instead of listing elements. For example, starting from the real numbers, we can consider all \(x\) such that \(\cos(x) > 1/2\); the collection of all such \(x\) is a set. Often we also fix a larger “universe” of objects, called the universal set and denoted by \(\Omega\), and only look at sets contained in this universe.

Given a universal set \(\Omega\) and a subset \(S \subseteq \Omega\), the complement of \(S\), denoted \(S^c\), is the set of all elements in \(\Omega\) that are not in \(S\):
\[
S^c = \{x \in \Omega : x \notin S\}.
\]
Taking complements twice brings us back to the original set: \((S^c)^c = S\). A special set is the empty set \(\emptyset\), which contains no elements. For example, the complement of \(\Omega\) is \(\emptyset\), since there is nothing outside the universal set.

If every element of a set \(S\) also belongs to a set \(T\), we say that \(S\) is a subset of \(T\) and write \(S \subseteq T\). This includes the case \(S = T\). Intuitively, \(S\) is “contained” inside \(T\).

\subsubsection{Union, intersection, and infinite collections}

Given two sets \(S\) and \(T\), their union \(S \cup T\) is the set of all elements that belong to at least one of them:
\[
x \in S \cup T \quad \text{iff} \quad x \in S \ \text{or} \ x \in T.
\]
Their intersection \(S \cap T\) is the set of elements that belong to both:
\[
x \in S \cap T \quad \text{iff} \quad x \in S \ \text{and} \ x \in T.
\]

These notions extend to infinite families of sets. If we have sets \(S_1, S_2, S_3, \dots\), their union is
\[
\bigcup_{n=1}^\infty S_n = \{x : x \in S_n \ \text{for at least one } n\},
\]
and their intersection is
\[
\bigcap_{n=1}^\infty S_n = \{x : x \in S_n \ \text{for all } n\}.
\]
Informally, the union collects elements that appear in at least one of the sets, while the intersection keeps only the elements common to all of them.

\subsubsection{Basic algebra of sets}

Set operations satisfy some familiar algebraic properties. Union and intersection are both commutative:
\[
S \cup T = T \cup S, \qquad S \cap T = T \cap S,
\]
and associative:
\[
(S \cup T) \cup U = S \cup (T \cup U), \qquad (S \cap T) \cap U = S \cap (T \cap U).
\]
Because of associativity, we usually drop parentheses when writing unions or intersections of three or more sets.

The universal set \(\Omega\) behaves like an identity: \(S \cup \Omega = \Omega\) and \(S \cap \Omega = S\). There are also distributive laws describing how union and intersection interact, for example
\[
S \cap (T \cup U) = (S \cap T) \cup (S \cap U),
\]
and a similar law with \(\cup\) and \(\cap\) swapped. A standard way to check such identities is to argue element-wise: assume \(x\) belongs to one side, unpack the definitions, and show that this forces \(x\) to belong to the other side as well. If each side is a subset of the other, then the two sets are equal.

\subsubsection{De Morgan's laws}

De Morgan's laws relate unions, intersections, and complements. For two sets \(S\) and \(T\), the first law says that the complement of their intersection is the union of their complements:
\[
(S \cap T)^c = S^c \cup T^c.
\]
Intuitively, if a point is not in both \(S\) and \(T\) at the same time, then it must be missing from at least one of them, so it lies in \(S^c\) or in \(T^c\). A formal proof rewrites the membership condition step by step: \(x \notin S \cap T\) means “\(x\) is not in \(S\) or not in \(T\),” which is the same as “\(x \in S^c\) or \(x \in T^c\),” that is, \(x \in S^c \cup T^c\).

The second law is symmetric and says that the complement of a union is the intersection of the complements:
\[
(S \cup T)^c = S^c \cap T^c.
\]
Here the idea is that if a point is not in \(S \cup T\), then it cannot be in \(S\) and cannot be in \(T\), so it must belong to both complements. Both identities extend to families of sets \(\{S_n\}\):
\[
\Big(\bigcap_{n} S_n\Big)^c = \bigcup_{n} S_n^c,
\qquad
\Big(\bigcup_{n} S_n\Big)^c = \bigcap_{n} S_n^c.
\]
These laws are very useful when we want to switch between unions and intersections, especially in probability.
\subsubsection{Sequences and their limits}

A sequence is a collection of elements indexed by the positive integers. We usually write it as \(\{a_i\}_{i=1}^\infty\), or simply \(a_i\), where each \(a_i\) belongs to some set \(S\). Often \(S = \mathbb{R}\) and we have a sequence of real numbers, but \(S\) could also be \(\mathbb{R}^n\) (a sequence of vectors) or some other set. Formally, a sequence is just a function \(f\) from the natural numbers to \(S\), and the \(i\)-th term is \(a_i = f(i)\).

We say that a sequence \(a_i\) converges to a real number \(a\) if, as \(i\) becomes large, the terms of the sequence get arbitrarily close to \(a\) and then stay close. Formally, \(a_i \to a\) as \(i \to \infty\) if for every \(\varepsilon > 0\) there exists an index \(i_0\) such that for all \(i \ge i_0\),
\[
|a_i - a| < \varepsilon.
\]
Geometrically, if you draw a horizontal band of width \(2\varepsilon\) around the limit \(a\), then from some point on all the sequence terms lie inside that band. The definition may look technical, but it just encodes the idea that the sequence eventually stays as close to \(a\) as we like.

Convergence behaves well under basic operations. If \(a_i \to a\) and \(b_i \to b\), then the sum sequence \(a_i + b_i\) converges to \(a + b\), and the product sequence \(a_i b_i\) converges to \(ab\). If \(g\) is a continuous function and \(a_i \to a\), then the sequence \(g(a_i)\) converges to \(g(a)\); for example, if \(a_i \to a\), then \(a_i^2 \to a^2\). These properties let us manipulate limits in a straightforward way.

\subsubsection{Practical convergence criteria}

In practice, it is helpful to have quick tests for convergence. One common situation is a monotone sequence. If \(a_i\) is increasing (or at least never decreases) and is bounded above, then it converges to some finite limit. If it keeps increasing without any bound, we say that it “converges to infinity” rather than to a finite number.

Another strategy is to bound the distance \(|a_i - a|\) from a candidate limit \(a\). If we can find an upper bound on \(|a_i - a|\) that goes to 0 as \(i\) grows, then \(a_i\) must converge to \(a\). A related idea is the sandwich (or squeeze) argument: if a sequence \(a_i\) is always between two other sequences \(b_i\) and \(c_i\), and both \(b_i\) and \(c_i\) converge to the same limit \(a\), then \(a_i\) also converges to \(a\). These tools will appear often, sometimes implicitly, when analyzing limits in the course.


\subsubsection{Infinite series}

Partiamo da una successione \(a_i\), per \(i = 1,2,\dots\), e proviamo a sommare tutti i termini:
\[
\sum_{i=1}^{\infty} a_i.
\]
Questa somma è definita come il limite delle somme parziali
\[
S_n = \sum_{i=1}^n a_i,
\]
se il limite di \(S_n\) esiste e è finito quando \(n \to \infty\). In tal caso diciamo che la serie converge; altrimenti non converge a un numero finito.

Se tutti i termini sono non negativi, le somme parziali \(S_n\) sono crescenti. Ogni successione monotona ha un limite (finito o \(+\infty\)), quindi la serie ha sempre un limite ben definito, anche se può essere infinita. Se i termini hanno segni diversi, il comportamento è più delicato: il limite può non esistere, oppure dipendere dall’ordine dei termini.

Per evitare problemi usiamo la convergenza assoluta. Se
\[
\sum_{i=1}^{\infty} |a_i|
\]
converge a un valore finito, allora anche la serie \(\sum a_i\) converge a un numero finito e la somma non cambia se riordiniamo i termini.

\subsubsection{The geometric series}

Un esempio centrale è la serie geometrica. Per un numero reale \(\alpha\), consideriamo
\[
\sum_{i=0}^{\infty} \alpha^i = 1 + \alpha + \alpha^2 + \alpha^3 + \dots
\]
Perché la serie converga, serve che \(\alpha^i \to 0\), cioè \(|\alpha| < 1\).

Per calcolare la somma guardiamo le somme parziali
\[
S_n = \sum_{i=0}^n \alpha^i
\]
e usiamo
\[
(1 - \alpha) S_n = 1 - \alpha^{n+1}.
\]
Se \(|\alpha| < 1\), allora \(\alpha^{n+1} \to 0\) e quindi
\[
(1 - \alpha) \sum_{i=0}^{\infty} \alpha^i = 1,
\]
da cui
\[
\sum_{i=0}^{\infty} \alpha^i = \frac{1}{1 - \alpha}, \qquad |\alpha| < 1.
\]

Un trucco più rapido: sia \(s = \sum_{i=0}^{\infty} \alpha^i\). Allora
\[
s = 1 + \sum_{i=1}^{\infty} \alpha^i = 1 + \alpha \sum_{i=0}^{\infty} \alpha^i = 1 + \alpha s,
\]
quindi \(s(1 - \alpha) = 1\) e \(s = 1/(1-\alpha)\). Questo ragionamento però presuppone che \(s\) esista, quindi va usato solo sapendo che la serie converge (ad esempio per \(|\alpha| < 1\)).
\subsubsection{Series with multiple indices and order of summation}

Sometimes we have series with two indices, for example \(a_{ij}\) with \(i,j = 1,2,\dots\). We can picture a grid of points \((i,j)\), each carrying a value \(a_{ij}\), and we want to sum all of them. Possibili modi: scegliere un ordine lineare dei punti e sommare in quell’ordine, oppure usare somme iterate:
\[
\sum_{i=1}^{\infty} \sum_{j=1}^{\infty} a_{ij}
\quad \text{oppure} \quad
\sum_{j=1}^{\infty} \sum_{i=1}^{\infty} a_{ij}.
\]
In teoria, ordini diversi possono dare risultati diversi.

La condizione che evita problemi è di nuovo la convergenza assoluta. Se
\[
\sum_{i=1}^{\infty} \sum_{j=1}^{\infty} |a_{ij}|
\]
è finita, allora la serie doppia è ben definita e qualsiasi ordine di somma (compreso scambiare \(i\) e \(j\)) dà lo stesso valore. Senza convergenza assoluta, si possono costruire esempi in cui un ordine dà 0 e un altro ordine dà 1, sfruttando cancellazioni tra termini positivi e negativi.

Possiamo anche avere somme doppie su insiemi ristretti, ad esempio \(\sum_{i=1}^{\infty} \sum_{j=1}^{i} a_{ij}\), che include solo i termini con \(1 \le j \le i\). Anche qui possiamo sommare prima in \(j\) e poi in \(i\), oppure prima in \(i \ge j\) e poi in \(j\). Con convergenza assoluta, queste diverse strategie portano allo stesso risultato.

\subsubsection{Countable and uncountable sets}

Molti spazi di probabilità sono infiniti, ma non tutti gli insiemi infiniti sono uguali. Alcuni sono discreti (countable), altri continui (uncountable). Un insieme è countable se i suoi elementi si possono mettere in corrispondenza biunivoca con i numeri naturali, cioè se possiamo elencarli in una sequenza: primo elemento, secondo, terzo, e così via.

Gli interi positivi sono ovviamente countable. Anche tutti gli interi lo sono: possiamo elencarli come \(0, 1, -1, 2, -2, 3, -3, \dots\), così ogni intero compare a una posizione finita. Lo stesso vale per tutte le coppie \((i,j)\) di interi positivi: possiamo percorrere la griglia a zig-zag e trasformare le coppie in una sequenza. L’idea si estende a triple, quadruple, ecc.

Un esempio meno ovvio è l’insieme dei razionali in \([0,1]\). Ogni razionale si può scrivere come \(p/q\) con \(p,q\) interi. Possiamo scorrere i denominatori \(q = 2,3,4,\dots\) e, per ogni \(q\), elencare tutte le frazioni \(p/q\) in \([0,1]\), saltando i duplicati. Alla fine otteniamo una sequenza che contiene tutti i razionali di \([0,1]\), quindi i razionali sono countable.

Un insieme è uncountable se non si può elencare in questo modo. Esempi tipici sono gli intervalli reali con lunghezza positiva, tutta la retta reale, oppure insiemi come il piano o lo spazio tridimensionale. Questi insiemi sono “troppo grandi” per essere messi in una semplice lista. La distinzione tra countable e uncountable è fondamentale per capire la differenza tra modelli di probabilità discreti e continui.[web:42][web:46]


\subsubsection{Cantor’s diagonal argument}

Cantor’s idea is to show that even a very small set of real numbers is already uncountable. Consider the numbers in \((0,1)\) whose decimal expansion uses only the digits 3 and 4, for example \(0.343434\ldots\). Each such number has an infinite decimal expansion made only of 3 and 4.

Supponiamo per assurdo che questo insieme sia countable. Allora potremmo elencare tutti i suoi elementi in una sequenza \(x_1, x_2, x_3, \dots\), scritti in base 10 usando solo 3 e 4. Costruiamo ora un nuovo numero \(y \in (0,1)\) così: il suo primo decimale è diverso dal primo decimale di \(x_1\) (se è 3 mettiamo 4, se è 4 mettiamo 3), il secondo decimale è diverso dal secondo decimale di \(x_2\), il terzo è diverso dal terzo di \(x_3\), e così via.

Per costruzione, \(y\) differisce da \(x_1\) nella prima cifra, da \(x_2\) nella seconda, da \(x_3\) nella terza, ecc. Quindi \(y\) non può coincidere con nessuno degli \(x_i\). Però \(y\) appartiene ancora al nostro insieme iniziale (sta in \((0,1)\) e usa solo 3 e 4), quindi la lista \(x_1, x_2, \dots\) non conteneva tutti gli elementi. Questo contraddice l’ipotesi che l’insieme fosse countable.

Dunque questo sottoinsieme dei reali è uncountable. Essendo un sottoinsieme di \(\mathbb{R}\), ne segue che anche l’insieme di tutti i numeri reali è uncountable. L’argomento diagonale mostra in modo molto chiaro che esistono infiniti “più grandi” di altri, e questo è alla base di molte idee nei modelli di probabilità continui.

